\section{Hindsight：保留、召回与反思的结构化智能体记忆}
\label{sec:hindsight}

Hindsight\upcite{latimer-etal-2026-hindsight}面向跨多轮对话和多个会话保持记忆的智能体。它把记忆组织为四类网络，并用retain、recall、reflect三种操作分别处理信息写入、证据检索和基于记忆的回答。

\subsection{研究问题与四类记忆}

作者将现有对话记忆的困难概括为三个相连的问题：长历史中怎样保留有用信息，查询时怎样跨会话找回相关信息，以及怎样在回答时避免只凭局部线索而产生全局不一致。单纯把片段放进向量库并取回 top-$k$，还可能混淆用户说过的事实与智能体形成的判断，也难以表达时间和实体之间的联系。

Hindsight 在每个 memory bank 中维护四类逻辑网络：
\begin{table}[!htbp]
\centering
\caption{Hindsight的四类记忆}
\label{tab:hindsight-types}
\begin{tabularx}{\linewidth}{@{}p{2.5cm}XX@{}}
\toprule
网络 & 保存内容 & 例子或作用 \\
\midrule
World & 关于外部环境的客观事实 & Alice 在 Google 的 Mountain View 团队工作。 \\
Experience & 智能体自身采取过的行动 & 智能体曾向 Alice 推荐徒步旅行。 \\
Observation & 根据多条事实整合出的实体摘要 & Alice 是专长于机器学习的软件工程师。 \\
Opinion & 带置信度、可随证据变化的主观判断 & “Python 适合数据科学”，并记录置信度。 \\
\bottomrule
\end{tabularx}
\end{table}
World与Experience记录事实或经历，Observation汇总实体证据，Opinion保存主观判断。这些类型区分了事实、概括与观点，便于检查和更新。记忆单元带有事件发生时间区间和系统获知时间，并通过时间、语义、实体和因果关系连接成图。

\subsection{三种操作：写入、检索与生成}

\textbf{Retain} 接收对话或文档，由模型抽取较粗粒度的叙事事实，再进行分类、实体识别与消歧、时间标注和关系建图。作者称每段对话通常抽取 2--5 条事实，以减少记忆切分方式对后续检索的影响。这里记录的时间至少要区分事件发生时间与智能体提及或获知该信息的时间。

\textbf{Recall} 接收查询和 token 预算，返回预算内排序后的记忆。它并行使用四类检索：HNSW 向量相似度、BM25 关键词匹配、沿实体／时间／因果连接进行图遍历，以及按事件时间范围过滤。不同通道的排序由 Reciprocal Rank Fusion 合并，再由 cross-encoder 重排。图遍历通常限制在 2--3 跳；时间解析先处理常见表达，规则无法处理时再调用轻量 T5 模型。token 预算让调用者限定送入回答模型的上下文规模。

\textbf{Reflect} 将查询、取回的记忆和可配置的行为档案交给生成流程。行为档案由怀疑程度、字面倾向和共情程度三个 1--5 的参数组成。Opinion 与事实分开存储，置信度在 $[0,1]$ 范围内；新证据支持某个判断时提高置信度，形成反证时降低置信度。后台整合还会保存支持性引文、证据数量和从 new 到 stale 的新鲜度趋势，使判断可以回溯到证据。

论文把 retain 和 recall 放在 TEMPR（Temporal Entity Memory Priming Retrieval）组件中，把 reflect 放在 CARA（Coherent Adaptive Reasoning Agents）组件中。端到端流程可概括为“对话进入四类记忆网络—多通道召回—依据证据生成并更新意见”。作者使用 PostgreSQL、pgvector 与全文检索索引作为存储基础，并报告：包含 10,000 条记忆的 bank，其 recall 检索延迟低于 200 ms，但该数字不包含回答模型的推理时间。

\subsection{实验设置与结果}

论文在LongMemEval与LoCoMo上评估长期对话记忆。LongMemEval的S设置包含500道问题，对话最长约150万tokens；LoCoMo使用多会话对话，每段最多35个会话。表\ref{tab:hindsight-results}列出不同模型配置下的回答准确率。

\begin{table}[!htbp]
\centering
\caption{Hindsight在两项基准上的回答准确率（\%）}
\label{tab:hindsight-results}
\begin{tabular}{lrrr}
\toprule
基准 & Hindsight 20B & Hindsight 120B & Hindsight Gemini-3 \\
\midrule
LongMemEval & 83.6 & 89.0 & 91.4 \\
LoCoMo & 83.2 & 85.7 & 89.6 \\
\bottomrule
\end{tabular}
\end{table}

在LongMemEval上，20B模型配合Hindsight达到83.6\%，120B模型为89.0\%，Gemini-3为91.4\%；在LoCoMo上，相应结果为83.2\%、85.7\%和89.6\%。20B配置在LongMemEval时间推理和多会话题上均为79.7\%，偏好题为66.7\%，反映出不同题型仍有差异。

作为参照，全文GPT-4o在LongMemEval上的分数为60.2\%，Supermemory-Gemini-3为85.2\%；LoCoMo中Zep、Memobase和Backboard分别为75.1\%、75.8\%和90.0\%。这些结果涉及不同基础模型和系统配置，适合描述各配置的表现；记忆机制的独立贡献则需在固定模型与预算下比较。

\subsection{方法局限与课设启发}

论文明确列出的局限包括：事实抽取、实体消歧和意见更新依赖语言模型，错误可能沿记忆图传播；评测只覆盖英文基准；Opinion 的主观质量和置信度校准还没有通过正式用户研究验证；时间解析可能漏掉含混或文化相关的表达。持续保存个人偏好与意见也带来隐私和偏见风险，论文指出系统没有内置 PII 检测或自动脱敏机制。

评测材料也有待完善：系统演示论文引用的完整评估条目仍带有arXiv编号占位符，且正文与附录中的部分LongMemEval分类基线数值不一致。复核组件效果时，需要补充完整实验配置和逐题结果。

Hindsight对课程设计的启发集中在记忆类型和检索组织：分别保存事实、经历、实体摘要与观点，利用时间和实体关系补充向量检索，并通过token预算控制回答上下文。可在固定模型、数据和预算下逐项比较检索通道，记录准确率、无关条目比例、延迟与token开销。

% 原文定位：PDF 第1--3页，印刷页275--277；PDF 第3--4页，§3.1、图1；PDF 第3--4页，§3.2--3.3；PDF 第3、5页，§3.2、§4.1；PDF 第4--5页，图1、§4.1--4.2；PDF 第7页及第11页，§6.1--6.3、表4；PDF 第8页，印刷页282；PDF 第7--9、11页，§6.2、表4及参考文献；PDF 第4页；PDF 第6页；PDF 第11页。
